\section{Introduction}
\label{sec:introduction}
%============high level big picture
This interdisciplinary project brings together formal and computational semantics with natural language processing (NLP). It examines how formal meaning representations can capture replacement self-repair (RSR) in natural language.

%=============intro starts with general definition
Speakers sometimes revise an utterance before they have finished it. This
process is called self-repair, a type of disfluency with three parts
\citep{shribergPreliminariesTheorySpeech1994,
leveltMonitoringSelfrepairSpeech1983}. The \emph{reparandum} is the abandoned
material. The \emph{interregnum} is the interval between the end of the reparandum and the onset of the coming repair. It may contain an editing
expression such as \textit{I mean} \citep{clarkPausesHesitationsPsycholinguistic2006}.
The \emph{repair} is the material with which the speaker resumes the utterance.
This project focuses on replacement self-repair (RSR) within a clause, in which the
repair replaces the reparandum and changes the intended meaning.

Figure~\ref{fig:siri-repair} illustrates an RSR: 
the speaker begins by asking Siri about the weather in Berlin, then replaces
Berlin with Dublin. In the utterance, \textit{in Berlin} is the reparandum,
\textit{I mean} is the editing expression occupying the interregnum, and
\textit{in Dublin} is the repair. Although
Dublin expresses the speaker's final intention, Siri returns the weather for
Berlin.

\begin{figure}[ht]
    \centering
    \includegraphics[width=0.45\textwidth]{viz/siri-b-i-d.png}
    \caption{Siri responds to the withdrawn location, Berlin, after the speaker
    replaces it with Dublin.}
    \label{fig:siri-repair}
\end{figure}

%=========quick LR showing previous method and limitations
Mainstream computational approaches commonly detect the reparandum and
editing expression for subsequent processing. One approach removes them while
retaining the repair, producing fluent text intended to express the speaker's
final intention \citep{johnsonTAGbasedNoisychannelModel2004}. Another retains
the spoken words and marks disfluent spans or tokens
\citep{futamiStreamingJointSpeech2023}. In the weather example, the first
approach retains only \textit{in Dublin}, while the second would preserve both cities but marks
\textit{in Berlin} as disfluent.

Explicit MRs remain useful in LLM-based systems because they make
model-interpreted meaning inspectable beyond surface text. They organise
entities, events and their relations into structures that can be examined
formally \citep{bosSequenceNotationCatching2023}. In recent work, explicit MRs
have been used to detect unsupported relations in generated database queries
\citep{reganMASSIVEMultilingualAbstract2024}, evaluate the contextual
appropriateness of dialogue responses
\citep{yangEmphasisingStructuredInformation2025} and construct
paraphrase-resistant semantic watermarks
\citep{yeSWANSemanticWatermarking2026}. These uses illustrate why explicit
semantic structure remains useful for verification, evaluation and subsequent
processing.

For applications using explicit MRs, self-repair raises the question of what
information should remain after a speaker revises an utterance.
The reparandum can continue to affect comprehension, as shown by
psycholinguistic evidence \citep{lauLingeringEffectsDisfluent2005}. It may also
remain available for later reference in dialogue
\citep{purverComputationalModelsMiscommunication2018}.

Returning to Figure~\ref{fig:siri-repair}, suppose the speaker receives the
correct response and follows with \textit{``Show me the airports of these two
cities''}, referring to both locations. Cleaned text would retain Dublin but
lose Berlin, leaving only one explicit antecedent for the plural expression
\textit{these two cities}. A transcript marked with disfluency labels would
preserve both names \citep{futamiStreamingJointSpeech2023}, but the labels would
not specify how Berlin's semantic role has been revised.

%==========turn back to my own work -- Claim
Thus, I argue that an MR for RSR should encode the semantic content of the
reparandum and repair, together with the structural configuration through which the
utterance is revised. In the running example, both cities should be linked to the weather
request, with an ordered replacement structure designating Berlin as the
withdrawn location and Dublin as the accepted one. This representation keeps
Berlin available for later reference while recording that Dublin has replaced
it in the original role. The same distinctions make a parser's output
inspectable: its errors can be traced to omitted content, incorrect status
assignments or incorrect links to shared content. The project has developed a
formal representation of within-clause RSR. The proposed thesis will validate
its coverage and train and evaluate semantic parsers that generate it.

%================WHY SBN
Simplified Box Notation (SBN)
\citep{bosQuantificationAnnotationDiscourse2021,bosSequenceNotationCatching2023}
is a suitable starting point because it combines a structured meaning
representation with a sequential format suitable for LLM generation. Its
structure can be extended to distinguish withdrawn from accepted content while
preserving their connections to the surrounding utterance. SBN is also used
for semantic annotation in the Parallel Meaning Bank (PMB), whose existing
sentence--SBN pairs provide a basis for synthesising RSR data
\citep{abzianidzeParallelMeaningBank2017}.

%=========================================RQs & Contribution

The proposed thesis will address two research questions.

\begin{enumerate}
  \item \textbf{Representation:} How should within-clause RSR
        be represented in SBN, and what are the coverage and
        limitations of the representation?
  \item \textbf{Parsing:} How well can LLM-based semantic parsers learn this
        representation from synthesised meaning-banked data, and what errors
        remain?
\end{enumerate}

This thesis aims to make three contributions.

\begin{enumerate}[itemsep=0.2em]
  \item A representation of within-clause RSR in SBN,
        together with an account of its coverage and limitations.
  \item A completed diagnostic study of how an LLM-based semantic parser
        trained on fluent sentence--SBN pairs processes RSR.
  \item A data-generation and training pipeline that constructs paired RSR
        utterances and repair-aware SBNs from meaning-banked data and trains
        LLM-based parsers to produce the repair-aware representation.
\end{enumerate}

%============Section Navigation

Section~\ref{sec:related-work} (\emph{Related Work}) reviews disfluency
processing and semantic accounts of repair. Section~\ref{sec:formalism}
(\emph{Formalism}) defines the repair-aware SBN. Section~\ref{sec:methodology}
(\emph{Methodology}) presents the diagnostic pilot, corpus construction and
model experiment. Section~\ref{sec:timeline} (\emph{Timeline}) records
completed work and remaining milestones.